\section{Related Work}
% 相关工作。
\label{sec:related}

%\arq{This is a really short related work section. One additional change, don't say "A prior workshop version", it deanonymizes us and we'll get rejected for it.  Just name the prior system.}

\texttt{sched\_ext}~\cite{schedext-docs}, \texttt{cache\_ext}~\cite{zussman2025cache_ext}, and XDP extend Linux subsystems via synchronous eBPF callbacks at fixed decision points;
% sched\_ext、cache\_ext 和 XDP 通过在固定决策点进行同步 eBPF 回调来扩展 Linux 子系统；
\sys generalizes this to asynchronous, cross-device effects on user-defined events (\S\ref{sec:design-interfaces}).
% 而 \sys 将其泛化为在用户自定义事件上产生异步、跨设备的效应 (\S\ref{sec:design-interfaces})。
Driver-level GPU resource managers~\cite{kato2011timegraph,kato2012gdev,kernelintercept,fan2025gpreempt,wang2024gcaps,coppock2025lithos} require unsafe OS kernel driver modifications and service interruptions to change policy.
% 驱动层 GPU 资源管理器~\cite{kato2011timegraph,kato2012gdev,kernelintercept,fan2025gpreempt,wang2024gcaps,coppock2025lithos} 需要不安全的 OS 内核驱动修改,并伴随服务中断才能更换策略。
User-space frameworks~\cite{ng2023paella,shen2025xsched,gim2025pie,chen2025ktransformers} offer programmability but lack cross-tenant visibility and privileged hardware control (\S\ref{sec:motivation-limitations}).
% 用户态框架~\cite{ng2023paella,shen2025xsched,gim2025pie,chen2025ktransformers} 提供可编程性,但缺乏跨租户可见性和特权硬件控制 (\S\ref{sec:motivation-limitations})。
\sys offers runtime-programmable policies without either limitation.
% \sys 在运行时提供可编程策略,且没有上述任何一种局限。
NVBit~\cite{villa2019nvbit} and Neutrino~\cite{neutrino} inject instrumentation into GPU binaries for device-side visibility, but lack safety guarantees and target offline profiling rather than online policy enforcement.
% NVBit~\cite{villa2019nvbit} 和 Neutrino~\cite{neutrino} 向 GPU 二进制中注入插桩以获得设备侧可见性,但缺乏安全保障,且面向离线性能分析而非在线策略执行。
A prior workshop paper~\cite{egpu} extends eBPF to GPU instrumentation but remains limited to read-only observability with high overhead (\S\ref{sec:eval}) due to per-thread scalar execution on SIMT hardware;
% 一篇先前的 workshop 论文~\cite{egpu} 将 eBPF 扩展到 GPU 插桩,但由于在 SIMT 硬件上按每线程标量执行,仍局限于只读可观测性且开销较高 (\S\ref{sec:eval});
\sys introduces SIMT-aware verification and warp-level execution for safe, low-overhead policy enforcement.
% \sys 引入 SIMT 感知验证和 warp 级执行,实现安全、低开销的策略执行。
